% Generated by tools/build_v6.py; edit the assembly rules, not this file.
\section{Representation and realization of the stopping flow}
\label{sec:flow-foundations}

This section justifies the normalized flow cone and robust program in
\eqref{eq:dual-flow-cone}--\eqref{eq:dual-primal}. Triangular encodings
of stopped coin experiments have direct antecedents in Lee--Valiant
\cite{leevaliant}; occupation-measure formulations of optimal stopping are
also classical \cite{cs}. The proofs below give the representation
and estimator realization needed for this particular robust program.

\subsection{The support function of the rarity class}

For bounded measurable $f:(0,1/2]\to\R$, put
\[
 Q_{\mu,\rho}(f)
 =\sup_{q\in\mathcal Q_\rho}\int q(d)f(d)\,\mu(dd).
\]

\begin{lemma}[A fractional upper-tail formula]\label{lem:ff-support}
The support function satisfies the attained identity
\begin{equation}\label{eq:ff-support}
 Q_{\mu,\rho}(f)
 =\min_{\ell\ge0}
       \left\{\rho\ell+\int(f-\ell)_+\,d\mu\right\}.
\end{equation}
The minimum can be restricted to
$[0,\|f_+\|_\infty]$. The supremum on the left is attained by selecting
the largest positive values of $f$, allowing a fractional selection
on the cutoff level set.
\end{lemma}

\begin{proof}
For every $0\le q\le1$ and $\ell\ge0$,
$qf\le q\ell+(f-\ell)_+$. Integration and $\int q\,d\mu\le\rho$
prove the upper bound by the right side.

If $\mu(f>0)\le\rho$, take $q=\1_{\{f>0\}}$ and $\ell=0$.
Otherwise choose a positive cutoff $\ell$ satisfying
\[
 \mu(f>\ell)\le\rho\le\mu(f\ge\ell).
\]
Such a cutoff exists by the monotonicity and one-sided continuity of
the tail distribution. Take $q=1$ on $\{f>\ell\}$, $q=0$ on
$\{f<\ell\}$, and choose a constant fraction in $[0,1]$ on
$\{f=\ell\}$ so that its total mass is $\rho$.
The resulting value is
\[
 \int_{\{f>\ell\}}f\,d\mu
   +\ell[\rho-\mu(f>\ell)]
 =\rho\ell+\int(f-\ell)_+\,d\mu.
\]
This proves equality and attainment, including atoms of $\mu$ or
atoms in the distribution of $f$. Cutoffs larger than
$\|f_+\|_\infty$ cannot improve on the endpoint.
\end{proof}

For an array $u=(u_{nk})$, define
\[
 L_u^\pm(d)
 =\sum_{n,k}u_{nk}\binom nk
          (1/2\pm d)^k(1/2\mp d)^{n-k},
 \qquad
 I_u^-=\int L_u^-\,d\mu,\qquad
 F_u=L_u^+-L_u^-.
\]
The root or final-depth terms are included only when the corresponding
array permits them. The convention $0^0=1$ covers the deterministic
Bernoulli endpoints. By the definition of $m(q)$,
\begin{equation}\label{eq:ff-robust-linear}
 \sup_{q\in\mathcal Q_\rho}\langle m(q),u\rangle
 =I_u^-+Q_{\mu,\rho}(F_u).
\end{equation}
In particular, the two bias suprema in \eqref{eq:dual-primal} are
\[
 c_0+I_z^-+Q_{\mu,\rho}(F_z-1),
 \qquad
 -c_0-I_z^-+Q_{\mu,\rho}(1-F_z).
\]
The moment and work suprema use \eqref{eq:ff-robust-linear} with
$u=t$ and $u=c$. Thus the robust program requires only integration
against the known magnitude law and a scalar cutoff optimization.

\subsection{Arbitrary seeds and ordered histories aggregate to the cone}

\begin{proposition}[Normalized prefix aggregation]
\label{prop:ff-aggregation}
Every capped one-unit policy with $1\le T\le N$, arbitrary independent
auxiliary randomization, and bounded terminal output has normalized
arrays in the cone \eqref{eq:dual-flow-cone} with root intensity one.
Its mean, raw second moment, and expected work under any $q$ are
respectively
\[
 \langle m(q),z\rangle,\qquad
 \langle m(q),t\rangle,\qquad
 \langle m(q),c\rangle.
\]
Multiplying the four arrays by $\lambda\ge0$ gives the corresponding
intensity-$\lambda$ representation. The stopping representation alone
requires no output or bounded-output assumption.
\end{proposition}

\begin{proof}
Place all auxiliary randomization in one seed $r$, drawn from a law
$\nu$ independent of the latent unit and its observations. Once the
seed is fixed, all decisions and outputs are deterministic functions
of the ordered observed prefix.

For each ordered binary prefix $x$ of length $n$, let $A_x$ be the set
of seeds for which the policy reaches $x$ and stops there; let $C_x$
be the set for which it reaches $x$ and continues. These sets are
disjoint. For $n\ge1$ their union is exactly $C_{x^-}$, where $x^-$
is the prefix with its last symbol removed. Write
\[
 \alpha_x=\nu(A_x),\quad
 \gamma_x=\nu(C_x),\quad
 \zeta_x=\int_{A_x}Y(r,x)\,\nu(dr),\quad
 \tau_x=\int_{A_x}Y(r,x)^2\,\nu(dr).
\]
At the root, $\gamma_\varnothing=1$ and $\alpha_\varnothing=0$;
at depth $N$, $\gamma_x=0$.

For $s=(n,k)$, sum these quantities over all prefixes with $k$ heads
and divide by $\binom nk$:
\[
 a_{nk}=\frac{\sum_{|x|=n,\,\#1(x)=k}\alpha_x}{\binom nk},
 \qquad
 c_{nk}=\frac{\sum_{|x|=n,\,\#1(x)=k}\gamma_x}{\binom nk},
\]
and analogously define $z_{nk}$ and $t_{nk}$ from $\zeta_x,\tau_x$.
Aggregating $\alpha_x+\gamma_x=\gamma_{x^-}$ gives
\[
 a_{nk}+c_{nk}
 =\frac{n-k}{n}c_{n-1,k}
    +\frac{k}{n}c_{n-1,k-1}.
\]
Cauchy--Schwarz on the disjoint union of the prefix-indexed
stopping-seed sets
gives $z_{nk}^2\le a_{nk}t_{nk}$. Hence the arrays belong to the
stated cone.

Conditional on Bernoulli mean $p$, every ordered prefix in this
count state has the same probability $p^k(1-p)^{n-k}$.
Summing its output integrals therefore gives the stated moment
identities after mixing over the latent state. Expected work is the
sum of the probabilities of continuing at all prefixes, including
the root: each continuation requests exactly one charged observation.
This proves the work identity and accounts for the first observation.

Multiplication by $\lambda$ preserves the linear flow identities and
the homogeneous perspective inequalities. The argument also applies
to outputs with finite first and second seed aggregates. For the
bounded outputs used in finite constructions, these are automatic.
\end{proof}

The normalized arrays encode probabilities over the independent seed
for fixed observation prefixes. They are not the unconditional
probabilities of reaching a count state. This distinction is the
reason the Pascal coefficients in \eqref{eq:dual-flow-cone} are
$(n-k)/n$ and $k/n$.

\subsection{Every feasible point is an observable estimator}

\begin{proposition}[Compound-Poisson realization]
\label{prop:ff-realization}
Every feasible point of \eqref{eq:dual-primal} constructs an observable
population estimator with uniform expected work at most $B$ and
worst RMS error at most $\sqrt{b^2+v}$. In particular,
\[
 R_{\mu,\rho}(B)\le P_N(B).
\]
\end{proposition}

\begin{proof}
Suppose first that $\lambda>0$. Draw a Poisson number of independent
genuine units with mean $\lambda$. Each unit first requests one
observation. At an observed state $(n,k)$, stop with probability
$a_{nk}/(a_{nk}+c_{nk})$ and continue otherwise. At a stopping state
with $a_{nk}>0$, return $z_{nk}/a_{nk}$. Unreachable states may have
arbitrary actions, and outputs at zero stopping flow can be set to
zero. The constraint $c_{Nk}=0$ enforces the cap.

For a fixed Bernoulli mean $p$, the expected number of units reaching
$(n,k)$ is
\[
 (a_{nk}+c_{nk})\binom nk p^k(1-p)^{n-k}.
\]
This follows by induction from the normalized Pascal recursion.
For example, a continuation at $(n-1,k)$ followed by a zero
contributes
$c_{n-1,k}\binom{n-1}{k}p^k(1-p)^{n-k}$, and the other parent
contributes the analogous term with $k-1$. Their sum is exactly the
displayed expression. Multiplying by the stopping or continuation
probability gives the same formula with $a_{nk}$ or $c_{nk}$ alone.

Return $c_0$ plus the sum of all terminal unit outputs. Independence
of the genuine units and the compound-Poisson moment identities give
\begin{align}
 \E_q\widehat\theta
 &=c_0+\langle m(q),z\rangle,\label{eq:ff-realized-mean}\\
 \Var_q(\widehat\theta)
 &=\left\langle m(q),\frac{z^2}{a}\right\rangle
 \le\langle m(q),t\rangle,\label{eq:ff-realized-variance}\\
 \E_qT&=\langle m(q),c\rangle.\label{eq:ff-realized-work}
\end{align}
Set $z_{nk}^2/a_{nk}=0$ where $a_{nk}=0$; the perspective constraint
then also gives $z_{nk}=0$. Every sum is finite, and each nonzero
terminal output is a finite real number. The root contribution to
\eqref{eq:ff-realized-work} is $\lambda$, since $m_{00}(q)=1$.
The work constraint therefore charges every first observation and
implies $\lambda\le B$.

The bias, moment, and work inequalities in the program prove the
claimed uniform guarantees. When $\lambda=0$, the flow recursion
forces $a=c=0$ and the perspective constraint forces $z=0$.
The deterministic estimator $c_0$ realizes the required bound;
any unused nonnegative moment slack is harmless.
\end{proof}

Allowing zero-observation marks would not improve the attainable
Poisson guarantee. Indeed, the event that a one-unit procedure stops
before its first observation depends only on its independent seed.
Poisson thinning separates those marks from the active units. Replace
their total random contribution by its deterministic expectation.
This preserves every population mean and the work, while removing
their nonnegative variance contribution. The remaining active
intensity, together with the adjusted offset $c_0$, has precisely
the root convention in \eqref{eq:dual-flow-cone}.

\subsection{Finite-support laws give finite second-order cone programs}

Suppose $\mu=\sum_{i=1}^{m}w_i\delta_{d_i}$ is supplied by its positive
masses and distinct positive magnitudes. For each of the four support
functions in the two bias, moment, and work constraints, introduce
one cutoff $\ell\ge0$ and nonnegative epigraph variables $r_i$ with
\[
 r_i\ge f(d_i)-\ell.
\]
The corresponding robust term is represented by
$\rho\ell+\sum_iw_ir_i$, by Lemma~\ref{lem:ff-support}.
All $f(d_i)$ are affine functions of the relevant flow or output
array. Thus these epigraph constraints are linear.

The perspective constraints are second-order cones:
\[
 z_{nk}^2\le a_{nk}t_{nk}
 \quad\Longleftrightarrow\quad
 \|(2z_{nk},a_{nk}-t_{nk})\|_2\le a_{nk}+t_{nk},
 \qquad a_{nk},t_{nk}\ge0.
\]
One additional cone represents an epigraph variable $u\ge b^2$,
for example
$\|(2b,u-1)\|_2\le u+1$, and the objective is $u+v$.

Consequently the finite-support version has $O(N^2+m)$ scalar
variables and cone or linear constraints, with a dense coefficient
description of size $O(mN^2)$ after the Bernoulli kernels are
evaluated. This is a dimension count at a specified cap, without a
uniform claim about the sufficient cap, numerical conditioning, or
bit complexity.
